\chapter{Building Space}
\label{stone:space}

Up to now, every construction in PhiBase has lived on a line. We measured lengths. We joined paths. We watched them grow. We multiplied them. We divided them. We even learned how to remember them. Everything happened along a single direction.

Buildings do not. Neither do bridges. Nor domes. Nor the beautiful stellated structures that began this journey. Sooner or later we must leave the line behind. We must learn how to build space.

\section*{The Obvious Answer}

At first the answer seems completely obvious. Space has three dimensions. Surely three numbers should be enough. That is exactly how most coordinate systems work. Every point is described by
\[
(x,y,z).
\]
Nothing could be simpler. For cubes. For boxes. For ordinary buildings.

But golden triangles are not ordinary. The moment we begin assembling them, something curious happens. The geometry keeps remembering how it was built. Three coordinates remember where we are. They do not remember how we arrived. That difference is more important than it first appears.

\section*{Thinking Like Builders}

Imagine two carpenters. The first measures the finished building. The second watches it being assembled. Both know the same structure. But they know it in different ways. One sees location. The other sees construction. PhiBase has always preferred construction. It remembers recipes instead of results. Perhaps space should do the same.

\section*{Six Families Of Beams}

Look carefully at a dodecahedron. Ignore the faces. Ignore the vertices. Instead, pay attention to the directions in which the beams naturally run. Something remarkable appears. There are six fundamental families. Every beam belongs to one of them.

Suppose we simply count. One beam from the first family. Three from the fourth. None from the others. Instead of writing
\[
(x,y,z),
\]
we write
\[
[a,b,c,d,e,f].
\]
Nothing has been measured. Everything has been built. The six numbers are not distances. They are construction counts.

\section*{Why Six?}

At first six numbers seem extravagant. Why not reduce them immediately to three? Because something precious would disappear.

Imagine describing a house. One description gives its street address. Another gives the complete set of blueprints. The address tells you where it is. The drawings tell you how to build it. The six-number description is like the drawings. It remembers the construction. Later we can always compute ordinary coordinates. But we never lose the recipe.

\section*{Edges Remember}

This is where the six-number system begins to feel natural. Take any edge of the construction. Walk from one end to the other. You will discover something surprising. Two of the three ordinary coordinates never change. Only one changes. The edge already knows which direction it belongs to. The six construction counts preserve that information automatically. The arithmetic is quietly remembering the geometry.

\section*{The Familiar Theme}

By now this pattern should feel familiar. Every PhiBase number remembers how it was constructed. Every path remembers how it was assembled. Now every point remembers how it was reached. The same idea has simply grown into another dimension. Nothing new has been invented. The construction has become richer.

\section*{Builder's Notebook}

Imagine assembling a small framework from six families of beams. Describe a point reached by one beam from the first family, two from the fourth, and one from the sixth. Do not ask where the point is. Ask instead, ``What did I build to get here?'' That question is the beginning of PhiBase space.

\section*{Looking Ahead}

Builders think in construction. Surveyors think in coordinates. Fortunately they are describing exactly the same world. The next stepping stone shows how to translate perfectly between those two languages. Once again, the way home turns out to be just as beautiful as the journey out.
